% ---------------------------------------------------------------------------
% Introducción (norma 3.12)
% ---------------------------------------------------------------------------
% NO se numera como capítulo (norma 2.6). El primer capítulo numerado es el que
% sigue inmediatamente a esta sección.
%
% La norma exige que contenga, en este orden: antecedentes; justificación e
% importancia del trabajo, con la hipótesis de manera implícita; planteamiento
% del problema; objetivo general; y objetivos específicos.
%
% Redactada en tercera persona. Los acrónimos se explican en su primera
% aparición dentro de esta sección, con su traducción cuando proceden de otro
% idioma (norma 2.4).

\seccionSinNumerar{Introducción}

Los equipos de operaciones de seguridad reciben, en la práctica totalidad de las organizaciones,
un volumen de alertas que supera su capacidad de análisis manual. La causa dominante no es la
falta de detección, sino el exceso de ella: los sistemas de correlación basados en reglas y
firmas ---el fundamento de la mayoría de las plataformas de detección y respuesta hoy
desplegadas--- generan un alto volumen de falsos positivos, y cada uno exige el mismo esfuerzo de
interpretación que una amenaza real, porque la alerta señala qué regla se disparó, no por qué eso
importa en el contexto concreto del activo afectado.

El presente trabajo, realizado como pasantía larga en la empresa Domotai C.A., aborda ese problema
mediante un motor de triaje que clasifica y prioriza alertas de seguridad de red, justifica cada
decisión con un modelo de lenguaje de gran escala (LLM, del inglés \emph{Large Language Model})
apoyado en generación aumentada por recuperación (RAG, del inglés \emph{Retrieval-Augmented
Generation}) ejecutada en un entorno local, y propone una acción de respuesta contenida sujeta a
validación humana antes de ejecutarse.

\subsection*{Antecedentes}

Las plataformas comerciales de detección y respuesta han evolucionado del \emph{endpoint} aislado
---EDR, del inglés \emph{Endpoint Detection and Response}--- hacia la correlación multicapa que
combina red, identidad, correo y nube en una sola plataforma, conocida como XDR (\emph{Extended
Detection and Response}); y hacia el servicio operado por un equipo humano que las administra las
veinticuatro horas, conocido como MDR (\emph{Managed Detection and Response}). Las evaluaciones
independientes de estas plataformas ---notablemente las de MITRE ATT\&CK Evaluations, que mide su
cobertura de técnicas adversarias contra la matriz pública MITRE ATT\&CK--- siguen registrando de
forma explícita el volumen de alertas que un analista debe triar como un criterio de comparación
entre proveedores, lo que confirma que la fatiga de falsos positivos persiste incluso en las
soluciones líderes del mercado.

En paralelo, los modelos de lenguaje de gran escala han empezado a incorporarse a la operación de
seguridad. Una primera línea de trabajo continuó el preentrenamiento de modelos de tipo
\emph{encoder} sobre corpus específicos de ciberseguridad ---informes de amenazas, bases de
vulnerabilidades, documentación técnica---, de lo que \textbf{SecureBERT} es la referencia original
\cite{aghaei2022securebert}. Una segunda línea publicó modelos generativos con preentrenamiento
continuado sobre el mismo tipo de corpus, de los que \textbf{Foundation-Sec-8B}, de Cisco
Foundation AI, es la referencia pública más reciente, con variantes que declaran explícitamente el
triaje y el resumen de casos de un centro de operaciones de seguridad (SOC, \emph{Security
Operations Center}) como caso de uso objetivo \cite{foundationsec8b}. Las plataformas comerciales
líderes ---Microsoft Security Copilot, CrowdStrike Charlotte AI y equivalentes--- han incorporado
asistentes basados en estos modelos, con un patrón consistente: redactan y resumen, no deciden; la
autoridad final permanece humana.

Complementando a los modelos generativos, la generación aumentada por recuperación demostró que
combinar un modelo de lenguaje con un sistema de recuperación sobre un corpus externo permite
anclar sus respuestas a documentos concretos y actualizar ese conocimiento sin reentrenar el modelo
\cite{lewis2020rag}. Esta técnica es la que este trabajo aplica para corregir la principal
limitación detectada en la justificación generada por un modelo pequeño: describe correctamente
qué campos cita, pero se equivoca en lo que esos campos significan.

\subsection*{Justificación e importancia}

Adoptar un modelo de lenguaje en el camino de una alerta de seguridad no es una decisión neutra.
El catálogo de riesgos de OWASP para aplicaciones basadas en LLM sitúa la inyección de instrucciones
como el riesgo número uno \cite{owaspllmtop10}, y en este dominio no es un riesgo hipotético: el
contenido de un registro de eventos ---el nombre de usuario probado, la cadena completa del
evento--- lo escribe quien está atacando. Si esos campos se concatenan al mensaje del modelo sin
distinguirlos de una instrucción, el atacante obtiene un canal directo hacia el sistema que decide
sobre su propio ataque. A esto se suma un segundo riesgo, más sutil: la alucinación con apariencia
de fundamento, en la que el modelo redacta una explicación plausible que cita un dato que la alerta
no contiene, y que un analista con prisa acepta precisamente porque suena a análisis correcto.

Estos riesgos, junto con la exigencia de que el análisis de una alerta no salga de la red del
cliente hacia un servicio externo, justifican dos decisiones de diseño que atraviesan todo el
trabajo: que los campos de una alerta se traten siempre como datos y nunca como instrucciones, y
que el modelo se ejecute en un entorno local, sin depender de una nube de terceros. La segunda
decisión responde además a una necesidad de negocio concreta de la empresa: Domotai C.A. ofrece
la automatización asistida por inteligencia artificial como una de sus líneas de servicio, y un
módulo que reduzca el ruido operativo de un cliente sin exponer sus datos ocupa un lugar dentro de
esa oferta que las plataformas comerciales, todas ellas dependientes de la nube del proveedor, no
cubren.

Por último, la restricción de continuidad de servicio es la que distingue a este trabajo de un
sistema de respuesta automatizada convencional: un triaje que acierta más pero que, por error,
interrumpe un servicio que el cliente presta a los suyos es peor que el sistema al que sustituye.
De ahí que la decisión de qué acción ejecutar nunca la tome el propio modelo de lenguaje, sino una
política determinista filtrada por un perfil de cliente configurable, con la validación de una
persona antes de cualquier acción que alcance un servicio.

\subsection*{Planteamiento del problema}

Diseñar este motor de triaje sobre el sistema de análisis de alertas que la empresa usa hoy exigía
partir de tres elementos que, en la práctica, no estaban disponibles al inicio de la pasantía: el
propio sistema de logs y correlación de un cliente real, un entorno donde ejercitarlo sin riesgo
para una operación productiva, y un conjunto de alertas con verdaderos y falsos positivos
identificados contra los que medir cualquier mejora. Sin una fuente de alertas no hay nada que
triar; sin un entorno controlado no hay dónde ejecutar ni verificar una acción de respuesta; y sin
un conjunto etiquetado no es posible calcular precisión, exhaustividad ni ninguna otra métrica de
clasificación.

El problema, formulado de manera acotada, es entonces el siguiente: dado que no se dispone de
acceso a un cliente ni a su módulo propietario de análisis de alertas, ¿es posible construir,
evaluar y demostrar en un entorno de laboratorio reproducible un motor de triaje que clasifique
alertas de acceso no autorizado a servicios de gestión expuestos, justifique cada decisión con
razonamiento verificable en vez de con la sola cita de una regla, y proponga una acción de
respuesta acotada sin comprometer la continuidad del servicio que protege, de forma que el
resultado sea representativo de lo que ese motor haría sobre la infraestructura real de un cliente
cuando esta esté disponible?

La respuesta de diseño fue modelar un cliente genérico en vez de esperar a uno concreto, y sustituir
el sistema de logs inexistente por Wazuh, una plataforma de detección basada en reglas desplegada
dentro de un laboratorio de red emulado con Containerlab \cite{containerlabdocs}, cuyo nivel de
regla cumple además el papel de método tradicional de referencia contra el que se mide el
prototipo.

\subsection*{Objetivo general}

Desarrollar un prototipo para la detección y el triaje de alertas de seguridad en entornos
empresariales, basado en modelos de lenguaje con un módulo de generación aumentada por
recuperación ejecutado sobre un modelo desplegado de forma local, de manera que permita clasificar
y priorizar alertas con razonamiento explicable, reduciendo los falsos positivos y apoyando la
labor de los analistas de seguridad.

\subsection*{Objetivos específicos}

\begin{enumerate}
  \item \textbf{Analizar y aprender a usar el módulo propietario.} Analizar la solución actual de
        la empresa para la recepción y el análisis de alertas, identificar sus limitaciones y las
        necesidades no cubiertas, y definir el caso de uso acotado sobre el que se construiría el
        prototipo.

  \item \textbf{Revisar el estado del arte.} Analizar las soluciones de detección y respuesta
        extendida y gestionada existentes en el mercado y las aplicaciones de modelos de lenguaje
        en el ámbito de la ciberseguridad, identificando sus ventajas, sus limitaciones y las
        necesidades no cubiertas por las herramientas actuales, para derivar de ello los
        requisitos funcionales y no funcionales del prototipo.

  \item \textbf{Configurar un entorno de pruebas.} Instalar y poner en funcionamiento un entorno
        con casos de prueba representativos, garantizando que los datos no salgan hacia servicios
        externos, y preparar un conjunto de alertas de prueba etiquetado ---verdaderos y falsos
        positivos--- que sirviera de referencia de evaluación.

  \item \textbf{Diseñar la arquitectura del sistema.} Definir el flujo de procesamiento, las
        reglas de decisión y sus umbrales, la forma de interacción con el entorno de pruebas, el
        formato del razonamiento explicable y los puntos de validación humana, así como las
        métricas de evaluación y el método tradicional de referencia contra el que se compararía
        el prototipo.

  \item \textbf{Implementar el prototipo.} Construir el caso de uso acotado, integrando la
        ingesta y normalización de alertas, la clasificación y priorización, la generación de
        justificación explicable y el flujo de validación humana en las decisiones críticas, con
        registro de trazas para auditoría.

  \item \textbf{Evaluar el prototipo.} Medir el desempeño del sistema frente al método tradicional
        de referencia definido en el objetivo anterior, y analizar los errores sistemáticos y las
        limitaciones del enfoque.

  \item \textbf{Elaborar la documentación técnica.} Generar la documentación y el informe técnico
        con la arquitectura, las decisiones de diseño, los resultados, las limitaciones y las
        líneas de trabajo futuro.
\end{enumerate}

\subsection*{Estructura del informe}

El documento se organiza en cuatro capítulos. El primero describe la empresa donde se desarrolló
la pasantía. El segundo reúne el marco teórico, recorriendo desde el problema del triaje de
alertas y las plataformas de detección y respuesta hasta los modelos de lenguaje, la generación
aumentada por recuperación, los agentes con herramientas, los riesgos de aplicar todo ello a este
dominio y las métricas con que se evalúa. El tercero expone la metodología seguida a lo largo de las seis fases del proyecto, y el
cuarto presenta los resultados obtenidos y su discusión. Cierran el documento las conclusiones y
recomendaciones, y las referencias consultadas.
